\section{Propositions}

This section states the central propositions concerning STOP as a diagnostic outcome within an enterprise modeled as a continuum (K\_EA). The propositions follow directly from the formal framing and the definition introduced in the preceding sections and are formulated at the level of inquiry rather than action \cite{churchman1971}. No proofs or illustrative examples are provided here.

\subsection*{Proposition 1 (Model-Relativity of STOP)}
A STOP outcome is always relative to a specified enterprise model. It is not an absolute property of the enterprise itself, but a statement about admissibility within the chosen modeling frame.

\subsection*{Proposition 2 (STOP as Admissibility Exhaustion)}
STOP occurs if and only if the continuation space defined by the enterprise model contains no admissible continuation state. In this case, continuation is not an analytically valid outcome under the given model, as admissibility is exhausted by the model’s internal constraints \cite{ashby1956}.

\subsection*{Proposition 3 (Independence from Analyst Competence)}
A STOP outcome does not imply analytical failure, insufficient effort, or lack of competence. It reflects a property of the modeled system under its constraints, not a deficiency of the analytical process.

\subsection*{Proposition 4 (Non-Equivalence of STOP and Inaction)}
STOP as a diagnostic outcome is not equivalent to organizational inaction. It expresses the non-existence of admissible continuation within the model and does not assert or prescribe any real-world decision or behavior. This separation between diagnostic representation and organizational action is consistent with established distinctions between model-level anticipation and real-world execution \cite{rosen1985}.

\subsection*{Proposition 5 (Symmetry of Analytical Outcomes)}
Within the formal framework, STOP and continuation are symmetric analytical outcomes. Each is determined solely by the structure of the admissible continuation space, and neither holds a privileged status a priori.

\subsection*{Proposition 6 (Model Revision Implication)}
The occurrence of STOP implies that any analytically valid continuation would require a modification of the model itself, rather than further exploration within the existing model. Such modification lies outside the scope of the STOP outcome as defined.

\subsection*{Proposition 7 (Diagnostic Informational Content)}
A STOP outcome carries positive diagnostic information. It identifies the boundaries of admissibility imposed by the current model and thereby delineates the limits of internally coherent evolution under that model.
